There is very little previous work that analyzes learning in the
setting of debate. Perhaps the two that are closest to our work
are~\citet{chen2023playing} and~\citet{buzaglo2025hidden}. Both model
debate as a repeated game between the two debaters with a fixed payoff
matrix. Since the space of arguments is too large to enumerate, the
debaters access this matrix through a best-response oracle, and both
papers give efficient no-regret strategies under this access. By
classic results, the time-averaged play then converges to an
equilibrium~\citep{freund1999adaptive,hart2000simple}. In this model,
the debaters never need to learn the verifier, because the oracle can
evaluate how the verifier would score any argument, and the
correctness of the equilibrium answer is assumed rather than analyzed.
In our model, the verifier is unknown, each round reveals a single
verdict on a new input, and we measure success by the number of
mistakes with respect to the ground truth.
